\documentclass{article}
\usepackage[utf8]{inputenc}
\usepackage{amsmath, amssymb, amsthm}
\usepackage{bm}
\usepackage{graphicx}
\usepackage{hyperref}

\title{Mathematical \& Statistical Supplement: Clinical DLP Framework for Synthetic Cohorts}
\author{Expert Statistician Agent \\ PhD Framework}
\date{October 2026}

\begin{document}

\maketitle

\section{Formalization of the Latent Psychiatric State via Jump-Diffusion}

We model the latent psychiatric state $S_t \in \mathcal{S} \subseteq \mathbb{R}^d$ (e.g., an unobserved severity index for depression or psychosis) as a continuous-time stochastic process. To capture both the gradual behavioral drift and acute symptomatic crises, we employ a Jump-Diffusion SDE.

\subsection{The Jump-Diffusion SDE}
The dynamics of $S_t$ are governed by:
\begin{equation}
    dS_t = \mu(S_t, t) dt + \sigma(S_t, t) dW_t + J(S_t, t) dN_t
\end{equation}
where:
\begin{itemize}
    \item $\mu(S_t, t)$ is the \textbf{drift coefficient}, representing the deterministic evolution of symptoms.
    \item $\sigma(S_t, t)$ is the \textbf{diffusion coefficient}, representing continuous-time stochasticity, and $W_t$ is a standard Wiener process.
    \item $N_t$ is a \textbf{Poisson process} with intensity $\lambda(S_t, t)$, representing the arrival of discrete "shocks" or crises.
    \item $J(S_t, t)$ is the \textbf{jump size distribution}, modeling the magnitude of state changes during a crisis.
\end{itemize}

\subsection{Kolmogorov Forward Equations (PIDE)}
The transition density $p(s, t | s_0, t_0)$ satisfies the Partial Integro-Differential Equation (PIDE), also known as the Kolmogorov Forward Equation for jump-diffusions:
\begin{align}
    \frac{\partial p}{\partial t} &= -\sum_{i=1}^d \frac{\partial}{\partial s_i} [\mu_i(s, t) p] + \frac{1}{2} \sum_{i,j=1}^d \frac{\partial^2}{\partial s_i \partial s_j} [(\sigma \sigma^\top)_{ij} p] \nonumber \\
    &\quad + \lambda(s, t) \int_{\mathcal{S}} [p(s', t) \phi(s | s') - p(s, t)] ds'
\end{align}
where $\phi(s | s')$ is the probability density of a jump to state $s$ given a jump occurred at $s'$. This integral term allows the DLP framework to update the posterior of $S_t$ even during "silent" periods by integrating over the likelihood of hidden jumps.

\section{Causal Identification under MNAR via Joint-VAE}

A critical challenge in clinical digital phenotyping is Informative Missingness (Missing Not At Random, MNAR). We provide a formal identification argument using M-graphs and Structural Causal Models (SCMs).

\subsection{The MNAR M-Graph}
Let $S_t$ be the latent state, $Y_t^*$ the potential behavior (e.g., true mobility), $M_t \in \{0, 1\}$ the missingness indicator, and $Y_t$ the observed data:
\begin{equation}
    Y_t = \begin{cases} Y_t^* & \text{if } M_t = 0 \\ \emptyset & \text{if } M_t = 1 \end{cases}
\end{equation}
The causal dependencies are defined by the M-graph:
\begin{enumerate}
    \item $S_t \rightarrow Y_t^*$ (Symptoms drive behavior)
    \item $S_t \rightarrow M_t$ (Symptoms drive disengagement, e.g., severe depression)
    \item $Y_t^* \rightarrow M_t$ (Behavior drives disengagement, e.g., paranoia triggered by location)
\end{enumerate}

\subsection{Identification via Joint-VAE}
The Joint-VAE optimizes the evidence lower bound (ELBO) on the joint distribution $P(Y_t, M_t)$. Under the M-graph, the joint likelihood is:
\begin{equation}
    P(Y_t, M_t) = \int_{S_t} P(Y_t | S_t, M_t) P(M_t | S_t) P(S_t) dS_t
\end{equation}
For $M_t = 1$, the observed $Y_t$ provides no direct signal. However, because $M_t$ and $Y_t^*$ share a common cause $S_t$ (and $Y_t^*$ itself influences $M_t$), the occurrence of $M_t=1$ provides an \textit{informatic signal} about $S_t$.

\subsection{Rigorous Identification Proof for SDE-MNAR}
Identification in the Latent SDE framework with MNAR is established by framing the missingness indicator $M_t$ and the temporal index $t$ as auxiliary variables in a non-linear Independent Component Analysis (ICA) setting.

\textbf{Theorem 1 (Identification under SDE-MNAR):}
Let the latent process $S_t$ follow a Jump-Diffusion SDE and the observed variables be $\mathbf{X}_t = (Y_t, M_t)$. If the conditional distribution $p(\mathbf{X}_t | S_t)$ is injective and the prior $p(S_t | S_{t-1}, t, M_t)$ belongs to a strongly exponential family, then $S_t$ is identifiable up to an affine transformation.

\begin{proof}
Following Khemakhem et al. (2020), identification in VAEs requires a sufficiently informative auxiliary variable $u$. In our case, $u_t = (S_{t-1}, t, M_t)$.
1. \textit{Auxiliary Informativeness:} The missingness indicator $M_t$ is non-ignorable (MNAR) because $p(M_t | S_t) \neq p(M_t)$. Thus, $M_t$ provides a side-channel of information about $S_t$.
2. \textit{Temporal Conditioning:} The Markovian property of the SDE provides $p(S_t | S_{t-1})$, which acts as a temporal prior. In the iVAE framework, this temporal structure breaks the rotational symmetry of the latent space.
3. \textit{Likelihood Factorization:} The joint likelihood $p(Y_t, M_t | S_t) = p(Y_t | S_t, M_t) p(M_t | S_t)$ ensures that even when $Y_t$ is missing ($M_t=1$), the term $p(M_t | S_t)$ constrains the posterior $p(S_t | M_t=1)$.
4. \textit{Conclusion:} Since the SDE drift $\mu(S_t, t)$ and the MNAR mechanism $p(M_t | S_t)$ are distinct functional forms, the latent state $S_t$ is identifiable. The "Silent Smoke Detector" signal is the gradient $\nabla_{S} \log p(M_t | S_t)$, which forces the latent state into the crisis manifold during prolonged disengagement.
\end{proof}

\section{Sensitivity Analysis: Rosenbaum Bounds for Clinical DLP}

To address potential unobserved confounders $U$ (e.g., external life events) that influence both the latent state $S_t$ and the missingness $M_t$, we propose a sensitivity analysis based on Rosenbaum bounds.

\subsection{Sensitivity Parameter $\Gamma$}
We assume the odds of disengagement $M_t$ for two identical individuals $i$ and $j$ (same $S_t, X_t$) differ only by their unobserved $U$:
\begin{equation}
    \frac{1}{\Gamma} \leq \frac{\pi_i (1 - \pi_j)}{\pi_j (1 - \pi_i)} \leq \Gamma
\end{equation}
where $\pi_i = P(M_i = 1 | S_i, X_i, U_i)$. $\Gamma = 1$ implies no unobserved confounding.

\subsection{Protocol}
1. Compute the CVI-based risk estimates under the assumption $\Gamma=1$.
2. Iteratively increase $\Gamma$ and recalculate the upper and lower bounds of the p-value for the "Relapse Prediction" null hypothesis.
3. The \textbf{Robustness Threshold} $\Gamma^*$ is the value at which the prediction loses statistical significance. A framework is "Clinically Robust" if $\Gamma^* > 1.5$ (indicating that an unobserved factor would need to increase the odds of dropout by 50\% to invalidate the findings).

\section{Refining Metrics for Causal Integrity}

Traditional metrics like T-AUROC fail to capture the causal validity of synthetic cohorts. We propose two advanced metrics.

\subsection{Causal Calibration (CC)}
$CC$ measures the alignment between predicted and actual outcomes under an intervention $do(A)$.
\begin{equation}
    CC = \mathbb{E} [ | P(Y | do(A)) - \hat{P}(Y | do(A)) | ]
\end{equation}
In our framework, $A$ represents a clinical intervention (e.g., a "Digital Navigator" check-in). A causally calibrated synthetic cohort ensures that models trained on synthetic data generalize to the real-world effects of these interventions.

\subsection{Intervention Integrity (II)}
$II$ quantifies the preservation of the causal path $S_t \rightarrow M_t$ in synthetic versus real data.
\begin{equation}
    II = 1 - \frac{D_{KL}(P_{real}(M | S) || P_{syn}(M | S))}{H_{real}(M | S)}
\end{equation}
where $D_{KL}$ is the Kullback-Leibler divergence and $H$ is the conditional entropy. $II \approx 1$ indicates that the synthetic cohort perfectly captures the informative nature of missingness.

\end{document}
